\documentclass[12pt, a4paper]{article}

\usepackage{ctex}
\usepackage{float, booktabs, parskip}
\usepackage{amsmath, amsthm, amssymb, mathrsfs, bm}
\usepackage{graphicx, subfigure, svg}
\usepackage{hyperref, endnotes}
\usepackage{listings}

\setlength{\parindent}{4em}
\setlength{\parskip}{1em}
\allowdisplaybreaks[4]


\newcommand{\di}{{\mathrm{d}}}


\title{{\Huge{\textbf{Homework 2 \\}}}}
\author{刘滨瑞 \quad 2021012579 \quad 未央-水木12}

\begin{document}

\maketitle

    \section{}

        临界底坡（Critical slope），即在某一流量条件下，正常水深$h_0$恰好等于临界水深$h_c$的底坡。

        当渠道的底坡大于临界底坡时，为陡坡（Steep slope），此时正常水深小于临界水深。
        当渠道的底坡小于临界底坡时，为缓坡（Mild slope），此时正常水深大于临界水深。

    \section{}

        比能沿程增加，即$\frac{\di E_s}{\di s} > 0$，这等价于$i > J_f$，进而等价于$h > h_0$。

        因此，在12条水面曲线中，$M_1$，$S_1$，$S_2$，$C_1$满足条件，是比能沿程增加的水面曲线。

    \section{}

    \subsection{}

        对于临界水深，有：$\frac{\alpha Q}{g} = \frac{A_c^3}{B_c^3} = \frac{(b h_c)^3}{b}$。

        取$\alpha = 1$，求得：$h_c = Q^{\frac{2}{3}} g^{-\frac{1}{3}} b^{-\frac{2}{3}} = 0.861 m$。

        对于临界底坡，有：$\frac{Q}{\sqrt{i_c}} = K_c = \frac{A_c}{n} R^{\frac{2}{3}} = (b h_c)^{\frac{5}{3}} (b + 2 h_c)^{-\frac{2}{3}} n^{-1}$。

        求得：$i_c = Q^2 (b + 2 h_c)^{\frac{4}{3}} n^2 (b h_c)^{-\frac{10}{3}} = 5.301 \times 10^{-3}$。

        注意到$i_c > i_1 > i_2$，因此两段渠道都是缓坡。

    \subsection{}

        \begin{figure}[htbp]
        \centering
        \subfigure{
\includegraphics[width=0.90\textwidth]{./assets/1.jpg}}
        \caption{水面线型式}
        \end{figure}

    \section{}

        从第一段的缓坡过渡到第二段的陡坡，水面会出现水跌，因此水面线会经过临界水深线，可以认为$h_c = 3 m$。

        取$\alpha = 1$，对于矩形渠道：$q = \sqrt{g h_c^3} = 16.27 m^2 s^{-1}$。

    \section{}

    \subsection{}

        \begin{figure}[htbp]
        \centering
        \subfigure{\includegraphics[width=0.60\textwidth]{./assets/2.png}}
        \caption{题5图}
        \end{figure}

        基本公式：
        $$\Delta s = \frac{E_{sd} - E_{su}}{i - \overline{J_f}}$$

        若认为缩断面后到水跃前的流动整体为渐变流，有：
        $$L = \frac{(h^\prime - h_{co}) + \frac{\alpha}{2g} (\frac{q^2}{h^{\prime 2}} - \frac{q^2}{h_{co}^2})}{i_1 - \overline{J_f}}$$

        假定$b \gg h$，流量模数：
        $$K = \frac{b h}{n} (\frac{b h}{b + 2 h})^{\frac{2}{3}} = \frac{b h^{\frac{5}{3}}}{n}$$

        损失为：
        $$J_f = \frac{(q b)^2}{K^2} = \frac{q^2 n^2}{h^{\frac{10}{3}}}$$

        对于$\overline{J_f}$，我们采用几何平均值的计算方法：
        $$\overline{J_f} = \sqrt{\frac{q^2 n^2}{h^{\prime \frac{10}{3}}} \frac{q^2 n^2}{h_{co}^{\frac{10}{3}}}} = \frac{(q n)^2}{(h^\prime h_{co})^{\frac{5}{3}}}$$

        代入数值，计算得：
        $$\overline{J_f} = 0.03960$$
        $$L = 77.99 m$$

    \subsection{}

        沿用5.1中的计算方式，将研究流域等分为$N$段进行计算。

        计算代码如下：

        \begin{lstlisting}[language=Python]
    from numpy import linspace, power, log2
    import matplotlib.pyplot as plt

    q = 5.61
    h_co = 0.62
    h_p = 1.06
    n = 0.025
    i1 = 0.01
    alpha = 1.0
    g = 9.8

    def J(h1, h2):
        return power((q * n), 2) \
                / power((h1 * h2), (5/3))

    def s(h1, h2):
        return ((h2 - h1) + \
                alpha * power(q, 2) / (2 * g) \
                * (power(h2, -2) - power(h1, -2))) \
                / (i1 - J(h2, h1))

    Ns = [1, 2, 3, 4, 6, 10,20, 40, \
            100, 200, 400, 1000, 4000, \
            20000, 100000] # Number of segments
    Ls = []
    for N in Ns:
        hs = linspace(h_co, h_p, N + 1)
        re0 = 0
        for i in range(N):
            re0 += s(hs[i], hs[i + 1])
        Ls.append(re0)

    print(f"L = {re0} m")

    plt.figure()
    plt.title("log(N)-L")
    plt.plot(log2(Ns), Ls, marker='o')
    plt.xlabel("log(N)")
    plt.ylabel("L")
    plt.savefig("Ls.png")
        \end{lstlisting}

        如图3所示，随着$N$的增长，$L$逐渐收敛，计算结果为：
        $$L = 81.32 m$$

        \begin{figure}[htbp]
        \centering
        \subfigure{\includegraphics[width=0.72\textwidth]{./assets/3.png}}
        \caption{算法收敛情况}
        \end{figure}

        图3表明，我们的数值算法计算平稳，收敛性良好。
        该算法等价于对$\frac{\di L}{\di s}$进行积分，因此其收敛结果应为$L$真实值。

        该方法得到的答案略大于整段计算得到的答案，相对误差为$4.1\%$，
        这说明我们在5.1中使用的渐变流假设是可取的，所造成的误差并不大。

        使用New Bing。对话记录参见图4。

        \begin{figure}[htbp]
        \centering
        \subfigure[Part 1]{\includegraphics[width=0.48\textwidth]{./assets/4.png}}
        \subfigure[Part 2]{\includegraphics[width=0.48\textwidth]{./assets/5.png}}
        \caption{与New Bing的对话记录}
        \end{figure}

        研究New Bing给出的代码，其解答思路基本正确，仅有一个小bug：25行的摩擦坡度前缺少一个负号。
        在修改后，代码运行结果为$L = 81.27 m$，与我们的答案一致，仅差$0.05 m$。

\end{document}
